% GENERATED FILE. Edit canonical lessons, then run npm run generate:books.
% canonical-source-hash: fnv1a64:82fe36f3ba926f66
% canonical-lessons: PT-C303-inundacao, PT-C303-inundar, PT-C303-afogamento, PT-C303-alerta, PT-C303-alertar

\chapter{Quando a água sobe}
\label{ch:pt-quando-a-agua-sobe}
\hlchaptermodality{303}

\begin{chapteropening}
\textbf{By the end of this chapter:} \emph{I can talk about floods and drowning, and about warnings and alerts.}

Complete the last lesson of chapter 303: I can talk about floods and drowning, and about warnings and alerts.
\end{chapteropening}

\section[a inundação]{a inundação — the flood}

\label{lesson:PT-C303-inundacao}



\begin{quote}
Before the new one: say the Portuguese for to collapse, then the Portuguese for the rubble.
\end{quote}

\subsection*{You'll want to know: a inundação}
\textbf{a inundação} — \textquotedblleft{}the flood\textquotedblright{}.

\begin{quote}
A inundação cobriu as ruas do centro.
\end{quote}

\begin{quote}
\emph{The flood covered the streets of the centre.}
\end{quote}

When a river overflows people also say \emph{a cheia}.

\begin{cousinweb}
\emph{inundação} is a learned borrowing from Latin \emph{inundātiō} \textquotedblleft{}a flooding\textquotedblright{} (EtymDB), built on \emph{unda} \textquotedblleft{}a wave\textquotedblright{}.
\end{cousinweb}

\begin{grammarlens}[title={What you've built}]
One more word for telling, explaining and describing at B1.
\end{grammarlens}

\subsection*{Your turn}
\begin{itemize}
\raggedright
  \item \emph{Say it:} \emph{a inundação}
  \item \emph{Say it:} \emph{a inundação}, once more
  \item \emph{Say it:} a sentence of your own with \emph{a inundação}
  \item \emph{Recall:} say the Portuguese for to collapse, then the Portuguese for the rubble, then say \emph{a inundação} again
\end{itemize}

\begin{tcolorbox}[breakable,colback=teal!4,colframe=teal!35!black,title={Before you move on}]
What does a inundação mean? (\textquotedblleft{}The flood\textquotedblright{}.) Say it once more.
\end{tcolorbox}
\section[inundar]{inundar — to flood}

\label{lesson:PT-C303-inundar}



\begin{quote}
Before the new one: say the Portuguese for the rubble, then the Portuguese for the flood.
\end{quote}

\subsection*{You'll want to know: inundar}
\textbf{inundar} — \textquotedblleft{}to flood\textquotedblright{}.

\begin{quote}
A chuva inundou a cave.
\end{quote}

\begin{quote}
\emph{The rain flooded the cellar.}
\end{quote}

Also of feelings or crowds: \emph{os turistas inundam a praça}.

\begin{cousinweb}
\emph{inundar} is a learned borrowing from Latin \emph{inundāre} \textquotedblleft{}to overflow\textquotedblright{}, built on \emph{unda} \textquotedblleft{}a wave\textquotedblright{} (Lewis \& Short).
\end{cousinweb}

\begin{grammarlens}[title={What you've built}]
One more word for telling, explaining and describing at B1.
\end{grammarlens}

\subsection*{Your turn}
\begin{itemize}
\raggedright
  \item \emph{Say it:} \emph{inundar}
  \item \emph{Say it:} \emph{inundar}, once more
  \item \emph{Say it:} a sentence of your own with \emph{inundar}
  \item \emph{Recall:} say the Portuguese for the rubble, then the Portuguese for the flood, then say \emph{inundar} again
\end{itemize}

\begin{tcolorbox}[breakable,colback=teal!4,colframe=teal!35!black,title={Before you move on}]
What does inundar mean? (\textquotedblleft{}To flood\textquotedblright{}.) Say it once more.
\end{tcolorbox}
\section[o afogamento]{o afogamento — the drowning}

\label{lesson:PT-C303-afogamento}



\begin{quote}
Before the new one: say the Portuguese for the flood, then the Portuguese for to flood.
\end{quote}

\subsection*{You'll want to know: o afogamento}
\textbf{o afogamento} — \textquotedblleft{}the drowning\textquotedblright{}.

\begin{quote}
Houve um afogamento na praia ontem.
\end{quote}

\begin{quote}
\emph{There was a drowning at the beach yesterday.}
\end{quote}

Built on \emph{afogar}, ‘to drown’.

\begin{cousinweb}
\emph{afogamento} is built on \emph{afogar} \textquotedblleft{}to drown\textquotedblright{}, which comes from Latin \emph{offōcāre} \textquotedblleft{}to choke\textquotedblright{}, built on \emph{faucēs} \textquotedblleft{}the throat\textquotedblright{} (Lewis \& Short).
\end{cousinweb}

\begin{grammarlens}[title={What you've built}]
One more word for telling, explaining and describing at B1.
\end{grammarlens}

\subsection*{Your turn}
\begin{itemize}
\raggedright
  \item \emph{Say it:} \emph{o afogamento}
  \item \emph{Say it:} \emph{o afogamento}, once more
  \item \emph{Say it:} a sentence of your own with \emph{o afogamento}
  \item \emph{Recall:} say the Portuguese for the flood, then the Portuguese for to flood, then say \emph{o afogamento} again
\end{itemize}

\begin{tcolorbox}[breakable,colback=teal!4,colframe=teal!35!black,title={Before you move on}]
What does o afogamento mean? (\textquotedblleft{}The drowning\textquotedblright{}.) Say it once more.
\end{tcolorbox}
\section[o alerta]{o alerta — the alert, the warning}

\label{lesson:PT-C303-alerta}



\begin{quote}
Before the new one: say the Portuguese for to flood, then the Portuguese for the drowning.
\end{quote}

\subsection*{You'll want to know: o alerta}
\textbf{o alerta} — \textquotedblleft{}the alert, the warning\textquotedblright{}.

\begin{quote}
O governo deu o alerta vermelho.
\end{quote}

\begin{quote}
\emph{The government issued a red alert.}
\end{quote}

\emph{Estar alerta} — to be on guard.

\begin{cousinweb}
\emph{alerta} was borrowed from Italian \emph{all’erta} \textquotedblleft{}on the lookout\textquotedblright{} (EtymDB), from \emph{erta} \textquotedblleft{}a height, a watchtower\textquotedblright{}.
\end{cousinweb}

\begin{grammarlens}[title={What you've built}]
One more word for telling, explaining and describing at B1.
\end{grammarlens}

\subsection*{Your turn}
\begin{itemize}
\raggedright
  \item \emph{Say it:} \emph{o alerta}
  \item \emph{Say it:} \emph{o alerta}, once more
  \item \emph{Say it:} a sentence of your own with \emph{o alerta}
  \item \emph{Recall:} say the Portuguese for to flood, then the Portuguese for the drowning, then say \emph{o alerta} again
\end{itemize}

\begin{tcolorbox}[breakable,colback=teal!4,colframe=teal!35!black,title={Before you move on}]
What does o alerta mean? (\textquotedblleft{}The alert, the warning\textquotedblright{}.) Say it once more.
\end{tcolorbox}
\section[alertar]{alertar — to warn, to alert}

\label{lesson:PT-C303-alertar}



\begin{quote}
Before the new one: say the Portuguese for the drowning, then the Portuguese for the alert.
\end{quote}

\subsection*{You'll want to know: alertar}
\textbf{alertar} — \textquotedblleft{}to warn, to alert\textquotedblright{}.

\begin{quote}
Os vizinhos alertaram os bombeiros.
\end{quote}

\begin{quote}
\emph{The neighbours alerted the fire brigade.}
\end{quote}

\emph{Alertar para um perigo} — to warn of a danger.

\begin{cousinweb}
\emph{alertar} is built on \emph{alerta}, from Italian \emph{all’erta} \textquotedblleft{}on the lookout\textquotedblright{} (EtymDB).
\end{cousinweb}

\begin{grammarlens}[title={What you've built}]
One more word for telling, explaining and describing at B1.
\end{grammarlens}

\subsection*{Your turn}
\begin{itemize}
\raggedright
  \item \emph{Say it:} \emph{alertar}
  \item \emph{Say it:} \emph{alertar}, once more
  \item \emph{Say it:} a sentence of your own with \emph{alertar}
  \item \emph{Recall:} say the Portuguese for the drowning, then the Portuguese for the alert, then say \emph{alertar} again
\end{itemize}

\begin{tcolorbox}[breakable,colback=teal!4,colframe=teal!35!black,title={Before you move on}]
What does alertar mean? (\textquotedblleft{}To warn, to alert\textquotedblright{}.) Say it once more.
\end{tcolorbox}
